\begin{textsample}{NEG Dim 4 – global\_north\_2024\_12 – Score 1.00 – global\_north\_2024\_12/119058566.txt}  \label{ex:f4_neg_131}
Children are among the victims of Israeli strikes on Gaza ' s besieged north and central areas, as humanitarian aid continues to be restricted.
Women and children were among the victims of the air raids, and were taken to the al-Ahli Arab Hospital Abood Abu Salama/Anadolu Agency Published On 19 Dec 202419 Dec 2024 Israel has carried out a series of attacks on northern and central Gaza that killed at least 40 people, including children, as a leading rights group accused it of committing"acts of genocide"by denying clean water to Palestinians.
Separate Israeli attacks on Jabalia in northern Gaza, where Palestinians have remained under a tight siege for more than two months, killed at least 16 people, including 10 members of the same family, on Thursday.
In Gaza City, an Israeli attack on a home in the Daraj neighbourhood killed at least four Palestinians, with another person killed in a strike on a group of civilians in Zeitoun neighbourhood.
Another 15 people were killed and dozens more wounded in Israeli air attacks on two schools -- the Dar al-Arqam and Shaaban al-Rayes schools -- sheltering displaced people in Gaza City.
Medics in Gaza told Anadolu news agency that most of the victims were women and children, while witnesses said the strikes caused huge destruction to the schools and nearby residential buildings.
Advertisement In central Gaza, a strike on Maghazi refugee camp killed at least four people.
There were fears the overall death toll would rise as many people were reported wounded in the Israeli attacks.
The Israeli military also issued forced evacuation threats to residents in Bureij refugee camp."The question is where people can go, as everywhere is overcrowded in the central areas of the strip,"said Al Jazeera ' s Tareq Abu Azzoum, reporting from nearby Deir el-Balah.
Wounded Palestinians including children were brought to the al-Ahli Arab Hospital following Israeli air raids on Gaza City Abood Abu Salama/Anadolu Agency At least 45,000 Palestinians have been killed and 107,000 others wounded in more than 14 months of attacks across Gaza, a densely populated territory that is at risk of famine and is experiencing emergency levels of hunger.
Israel began its ferocious military campaign after at least 1,139 people were killed during an attack led by Palestinian group Hamas in Israel on October 7, 2023, with more than 200 others taken captive.
Most of Gaza ' s 2.3 million inhabitants have been displaced, many of them multiple times, while Israel ' s intense bombardment has left much of the territory in ruins.
Referring to a report by Human Rights Watch ( HRW ) on Thursday that accused Israel of using water as a weapon of war in Gaza, Abu Azzoum said searching for water in the besieged and bombarded territory was"a daily struggle for survival".
In its 179-page report, HRW detailed how Israeli authorities have cut off and later restricted piped water to Gaza, rendered most of its water and sanitation infrastructure useless by cutting electricity and restricting fuel, deliberately destroyed and damaged water and sanitation infrastructure and water repair materials ; and blocked the entry of critical water supplies.
Advertisement Israel ' s Ministry of Foreign Affairs denied the accusations, which it called"lies", and alleged that the organisation is promoting"anti-Israel propaganda".
A separate report by Doctors Without Borders ( MSF ) published on Thursday found"clear signs of ethnic cleansing"in Gaza, particularly in the northern part of the strip."Our firsthand observations of the medical and humanitarian catastrophe inflicted on Gaza are consistent with the descriptions provided by an increasing number of legal experts and organisations concluding that genocide is taking place in Gaza,"the organisation said.

% matched lemmas: 
\end{textsample}
